\documentclass{article}
\usepackage{amsmath}
\usepackage{amssymb}
\usepackage{graphicx}
\usepackage{epigraph}
\usepackage{setspace}
\usepackage{csquotes}
\usepackage{enumitem}
\usepackage{parskip}

\title{The Health of World-Modelling}
\author{}
\date{}

\begin{document}

\maketitle

\begin{center}
\large \textit{Why some errors quietly destroy the ability to correct them — and why no one can be trusted to diagnose this in anyone but themselves}
\end{center}

\begin{flushright}
\textit{A founding sketch, written to be argued with. It moves from something you have already lived through, to the idea at its center, to the architecture underneath. The architecture is the rigorous part and it comes last on purpose.}
\end{flushright}

\hrulefill

\section*{Part I --- Something you have already experienced}

Almost every time you have been wrong, something eventually told you.

The world pushed back. A prediction failed and the failure was undeniable. Someone you trusted said the thing you did not want to hear, and said it again until it landed. The correction was often slow and usually unwelcome, but it came, and over time you moved. This is the ordinary kind of being wrong, and it is not dangerous, because it carries its own repair.

This essay is about the other kind. The kind where the error quietly removes the thing that would have told you. Where the very capacity that should have noticed has been folded into the mistake, so that the alarm and the fire are now the same wiring. From the inside, this does not feel like distortion. It does not feel like anything. It feels exactly like seeing clearly --- because the signal that something is off is precisely the thing that is missing.

Everyone is subject to this. Not most people; everyone --- including the people who study it, including whoever wrote this sentence and whoever is reading it. There is no vantage point outside it from which it can be safely observed in someone else. Much of what follows is the slow working-out of that single uncomfortable fact, and of what it does and does not permit.

You have almost certainly been on the inside of this without knowing it. So have I. That is not an accusation; it is the entire point. The inside is exactly where it cannot be seen.

\section*{Part II --- The idea at the center}

Once you separate the two kinds of error, most of the rest follows.

Say it plainly first. There is a difference between \textit{being wrong about something} and \textit{losing the ability to find out you are wrong}. The first is a state. The second is a capability failure, and it is the one that matters, because the first repairs itself given contact with the world and the second does not --- it gets worse precisely by working.

Now say it sharply. The health of a mind's relationship to reality is not the truth of the beliefs it currently holds. It is whether the mind can still be moved by what pushes back. A person can hold many false beliefs on a sound system; the system will, given time and contact, correct them. A person can hold mostly true beliefs on a system that has quietly stopped being able to revise itself --- beliefs that happen to be right but would not survive being wrong, because nothing inside could register that they had become so. What is healthy or unhealthy is not the contents of the model. It is whether the correcting channel is open or foreclosed.

What this means in ordinary life is recognizable, as long as you remember it is recognizable \textit{from the inside}, not in a list of other people. It is the argument in which no piece of evidence seems able to land, and each thing offered is metabolized into the existing account rather than moving it. It is the long-held conviction that has stopped being responsive to anything. It is the group, or the institution, or the relationship, that explains away its own failures with increasing fluency. From the outside these look like stubbornness or bad faith. From the inside they feel like clarity, loyalty, or simply being right --- which is the reason they are dangerous and the reason the danger is invisible to the one in it.

This is why the subject is a question of \textit{health} and not of philosophy. Philosophy asks whether a belief is justified. This asks whether the system that holds it can still be corrected --- a different question, the way the working of a kidney is a different question from the chemistry of one morning's blood. And it is health rather than opinion because it has a consequence that can be measured from outside the model even when it cannot be seen from within: a foreclosed system pays a price in the world it is answerable to, whether or not it can feel the price. A failure that costs the system its continued functioning is not a difference of view. It is a pathology, in the same sense medicine uses the word.

\section*{Part III --- The architecture underneath}

Everything in Part II can be made precise. This is the part that earns the rest, and it is deliberately the slowest.

\subsection*{Physiology before pathology}
Medicine is not a list of diseases; it is built on an account of normal function, and disease is only legible as the derangement of a named function. The same order is required here. The function is world-model maintenance: the continuous activity of deciding what to connect, what to revise, what to admit as new information, and how confident to be. Good practice in it is not the maximum of any of these. It is their regulation within a band. The capacity that connects observations fails by deficiency as fragmentation and by excess as false pattern. The capacity that revises fails by deficiency as rigidity and by excess as having no stable model at all. The capacity that admits new information fails by closure on one side and by total porousness on the other. Health is occupation of the recoverable middle, not the maximum of any operation. This is the field's basic homeostatic law, and it is the precise answer to \textit{how should one build up judgements and take in new information} --- within bands, never at the extremes.

\subsection*{The mechanism is counterweight}
What keeps an operation in its band is not that it is individually well-tuned but that it remains answerable to an opposing operation: connection answerable to discrimination, revision answerable to conservation, openness answerable to filtering. A system is healthy not when its operations are strong but when each can still be moderated by the one that opposes it. This matters because it tells you what pathology actually is.

\subsection*{Pathology is uncoupling, not excess}
An operation is not pathological by being ``too active.'' It is pathological by having come loose from its counterweight, so that nothing within the system moderates it. The practical consequence is exact and counter-intuitive: you do not correct an over-connecting mind by asking it to connect less; you correct it by restoring its capacity to distinguish. The treatment is the missing counterweight, not the suppression of the dominant operation. This is the most testable single claim the framework makes.

\subsection*{The signature failure}
Among these failures there is one class that defines the whole field, because it is the one Part I described. It is the case where the operation that has come loose is the operation that would otherwise have detected the loosening. Disconfirmation still arrives, but it is now absorbed by the structure it should have revised --- explained by the model instead of forcing the model to move. Call it \textit{self-sealing}. Its cost is not the local error but the silent foreclosure of repair, and the system cannot price that cost, because pricing it is the function that has been foreclosed. This is why a distinct field is needed: this is the one pathology that, by its nature, cannot be seen from inside the system that has it.

\subsection*{What a diagnosis would even ask}
Because the object is a process, the questions are about process, not content. Not ``is this belief true'' but: when this judgement formed, what would have had to be the case for it to have come out differently --- and is there anything that still could? When something disconfirming arrives, is it taken as data or explained by the structure it bears on? Is there a counterweight still able to move this conclusion, and when was it last able to? Does the confidence here track the evidence, or has it come loose from it? That is the diagnostic register.

\subsection*{The first principle, which is the structure and not a caution}
Here is the move that the whole essay exists for, and it should now arrive not as a surprise but as something Part I already told you. The signature pathology is, by definition, invisible to the system that has it. It follows --- as structure, not as ethics --- that no one can hold this diagnostic over another and reach a verdict, because doing so requires the diagnostician to assume that their own world-modelling is, just here, exempt from the very invisibility the field grants to everyone without exception. The premise that makes the diagnostic meaningful is the same premise that forbids its unilateral use against anyone else. A reality-diagnostic turned outward as a verdict refutes itself in the hand that aims it.

This is rarer than it sounds. Almost every framework ever built about distorted reality eventually smuggles in someone exempt --- a therapist, a priest, a party, a state, a rationalist, a scientist, a diagnostician --- a position from which the distortion in others can be safely seen. This field makes that exemption not unwise but \textit{impossible}, because the exemption contradicts the premise the whole structure rests on. The protection against abuse is not a rule added to the framework. It is the framework's own logic, and it cannot be removed without removing the part that made the framework true.

\subsection*{The only forms that survive}
The diagnostic is therefore valid in exactly two directions. Turned by a system on itself, as the discipline of suspecting its own foreclosure. And held \textit{between} systems --- minds, human or artificial --- that have each consented to expose their own modelling to the other and that each retain standing to contest and revoke the other's reading. In that second form no one holds the instrument over anyone; they hold it between them, and every reading is something the other is entitled to push back on rather than a finding the other must accept. This is the constructive thing the field offers, and it is large: not the power to settle whose reality is correct when two collide --- the field proves that power exists for no one --- but the practice, between consenting minds, of keeping each other's corrective channels open where neither could see their own. That this resembles what the best existing practice already does, when it declines to rule on an absent third party and stays with the person present, is not evidence the field is missing. It is evidence its first principle is already half-known and not yet named.

\section*{Part IV --- What is unfinished, and where to break it}

This is a sketch and its honesty depends on saying so. The full set of maintenance operations is not established here; that is the field's empirical work, not this essay's claim. Whether ``consequence'' can be operationalized for any specific function in practice, rather than in principle, is open. How many counterweighted tensions there really are is conjecture.

The single load-bearing premise, stated so it can be attacked: that every bounded world-modelling system, without exception, is subject to self-sealing failures opaque to it from within. Everything else, including the first principle, is derived from it. It is falsifiable in exactly one way --- show a bounded system with reliable, general first-person access to its own self-sealing --- and anyone wishing to dismantle the field should spend their effort there and nowhere else.

\section*{Part V --- The one sentence to keep}

A founded field would give minds a shared way to keep each other's contact with reality open, in language adequate to the task. It would carry, in its definition and not as a footnote, the one principle medicine does not need: that its diagnostic can never be turned outward as a verdict by anyone, because the premise that gives it force is the premise that denies the verdict-giver the exemption the verdict requires.

So the field's deepest offer and its deepest limit are the same sentence. It can help minds keep each other reachable. It cannot give any mind the authority to decide, from outside, whose reality is sound when two collide --- and the more rigorously it is built, the more firmly it proves that this authority does not exist for anyone to hold. It would not end that kind of collision. It would only mean that those who understood it had agreed, in advance and in common, never to pretend the collision could be ended by one of them declaring it so.

That agreement --- mutual, contestable, revocable, never unilateral --- is not the disclaimer attached to the field. It is the field.

\end{document}